\section{View lines, deflection, and timing}
\label{sec:rays}
\subsection{Null propagation in the specified metric}
For numerical ray tracing define
\begin{equation}\label{eq:optical}
 N=\sqrt{1+2\ph},\qquad B=\sqrt{1-2\ph},\qquad
 v=\frac{N}{B}=1+2\ph+\order(\epsilon^2).
\end{equation}
The symbol $v$ is a dimensionless coordinate speed per conformal time on
the unit sphere. The speed measured locally by an orthonormal observer
remains $c$. The code treats the displayed truncated metric as its
numerical metric. Taking square roots and integrating its trajectories
retains some terms of order $\epsilon^2$ and higher, but does not promote
it to an exact solution of the nonlinear Einstein equations. All physical
claims are restricted to first order.

Set $s=\eta_o-\eta$. Spherical symmetry about the mass confines each
spatial ray to a totally geodesic two-sphere in $\Sthree$. Represent this
two-sphere by $\bm n\in\mathbb R^3$, with
$\bm n^2=1$ and tangent momentum $\bm p\cdot\bm n=0$. The mass lies
at $\bm n_L=(1,0,0)$, so $\mu=\bm n\cdot\bm n_L$.
The reduced optical Hamiltonian is $\mathcal K=v|\bm p|$. In the
past-directed parameter and with momentum pointing along the past
spatial trajectory, its equations are
\begin{align}
 \frac{\dd\bm n}{\dd s}&=v\frac{\bm p}{|\bm p|},\label{eq:ray-n}\\
 \frac{\dd\bm p}{\dd s}&=-|\bm p|\nabla_S v-v|\bm p|\bm n,
 \label{eq:ray-p}\\
 \nabla_Sv&=\frac{2v}{1-4\ph^2}\,
 \epsilon g F_{,\mu}(\bm n_L-\mu\bm n).
 \label{eq:gradient}
\end{align}
The last term in Eq.~\eqref{eq:ray-p} preserves tangency to the unit
sphere. These equations preserve $\bm n^2=1$, $\bm n\cdot\bm p=0$,
and the angular momentum component about the lens axis. Conformal time
is not an affine parameter.

The physical ray is displayed using the \emph{background-time} map
$\bm X_{
\mathrm{view}}=A(\eta)\bm n(\eta)$, which keeps radius exclusively as
the time-layer label. Figure~\ref{fig:nested} uses this map.
Equation~\eqref{eq:deformed-embedding}, used separately for spatial
deformation, has a different purpose. The two maps are explicitly
distinguished instead of silently assigning both meanings to one radius.

\subsection{A deflection with fixed endpoints}
Comparing arbitrary curves at equal coordinate parameters is insufficient
to define a lensing prediction. We fix a coordinate-stationary source
worldline, an observer worldline, and a reception event. For the reference
FLRW ray, choose a launch angle $\alpha_0$ in the observer's tangent frame
and a source distance $s_0$. Its source position is
\begin{equation}
 \bm n_s=\cos s_0\,\bm n_o+\sin s_0\,\bm e_o(\alpha_0).
\end{equation}
For the perturbed ray, a two-variable shooting calculation adjusts
$(\alpha,s)$ until the same spatial source is reached. The emission time
is allowed to change because the source is a worldline, not a fixed
emission event. The predicted angular displacement is
\begin{equation}\label{eq:angular-shift}
 \delta\alpha=\alpha-\alpha_0.
\end{equation}
The angular coordinate in the observer's tangent plane equals the
orthonormal-frame angle because the spatial metric there is conformally
isotropic. Consequently Eq.~\eqref{eq:angular-shift} is an astrometric
model-to-reference prediction for the stipulated source matching. It is
not inferred from the apparent winding of a polar drawing.

The sources and observer in this experiment are test worldlines held at
fixed coordinates. They are not assumed to follow the dust velocity
field. Using dust tracers would require evolving their positions and
including their peculiar-velocity contributions.

An independent first-order check follows by projecting the perturbed ray
onto the parallel-transported transverse unit vector $\bm b$ along its
reference great circle. The transverse displacement obeys
$y''+y=-2\bm b\cdot\nabla_S\ph$, with the derivatives here taken
with respect to $s$. Setting $y(0)=y(s_0)=0$ gives
\begin{equation}\label{eq:born-angle}
 \delta\alpha_{\mathrm{first}}=\frac{2}{\sin s_0}
 \int_0^{s_0}\sin(s_0-s)\,
 \bm b\cdot\nabla_S\ph(\eta_o-s,\bm n_0(s))\,\dd s.
\end{equation}
The test remains away from $s_0=\pi$, where antipodal focusing would
invalidate this simple boundary-value comparison.

\subsection{An independent radial timing calculation}
A ray through the center follows a meridian with no transverse deflection.
Let $\ell$ be its unwrapped spatial arc coordinate. The null condition
gives $\dd\ell/\dd s=v$. For a fixed source arc distance $s_0$, expand
the arrival condition $\ell(s)=s_0$ to first order about the background
path $\bm n_0(\ell)$:
\begin{equation}\label{eq:born-delay}
 \Delta s=-2\int_0^{s_0}
 \ph(\eta_o-\ell,\bm n_0(\ell))\,\dd\ell+\order(\epsilon^2).
\end{equation}
This is an additional conformal lookback at fixed reception, relative to
the matched FLRW source. A negative potential along most of the path
gives $\Delta s>0$: the emission must occur earlier. The corresponding
background cosmic-time difference is computed with Eq.~\eqref{eq:dust}.
Raw conformal-time differences depend on the specified coordinate and
background identification. An operational clock delay would additionally
specify the emission/reception clock protocol.

\subsection{Frequency along a view line}
For the coordinate-stationary normal observers, the photon energy is
proportional to $|\bm p|/(AB)$. Thus
\begin{equation}\label{eq:redshift}
 1+z=\frac{|\bm p_e|}{|\bm p_o|}\frac{A_oB_o}{A_eB_e}.
\end{equation}
The normalization of $\bm p$ cancels. At $\epsilon=0$, this reduces to
$1+z=A_o/A_e$, checked numerically. For moving sources or receivers,
the appropriate four-velocities must instead be used in
$(k_\mu u^\mu)_e/(k_\mu u^\mu)_o$. The present calculation includes
the metric's gravitational and expansion effects for the stated normal
observers, and omits dust Doppler terms.
